\documentclass[a4paper]{article}

\usepackage[english]{babel}
\usepackage[utf8]{inputenc}
\usepackage{amsmath}
\usepackage{graphicx}
\usepackage[colorinlistoftodos]{todonotes}
\usepackage{hyperref}

\title{CHATBOTS AND CUSTOMER SATISFACTION: EVALUATING THEIR
IMPACT ON CUSTOMER EXPERIENCE AND BUSINESS BENEFITS}

\author{\makebox[\textwidth]{Ibrahim Abdo (231018), Jason van Hamond (232567), Radna Puriel (234908), Maria Salop (230574)}}



\date{October 14, 2024}

\begin{document}
\maketitle
\section{Executive summary}

\subsection{Problem Statement}
The rapid adoption of digital transformation has become essential for small and medium enterprises (SMEs) to remain competitive, with technologies like Artificial Intelligence (AI) playing a pivotal role in this shift. Digiwerkplaats, an organisation supporting SMEs in their digital transformation, aims to help businesses navigate the opportunities and challenges posed by AI. This research will explore the perception and impact of chatbots among consumers and business performances. By evaluating the risks and opportunities, this study will provide insights to guide stakeholders toward optimal AI implementation strategies.

\subsection{Project Objective}
The project objective of this research is to explore the impact of chatbot functionalities on customer satisfaction, trust, and overall customer experience and provide valuable insights for businesses. Specifically, the paper will examine: 

\begin{enumerate}
    \item \textit{The Customer Perception of Chatbot Reliability.}
    The research will focus on how the perceived reliability of chatbot-provided information, such as accuracy and clarity, influences customer trust to use the chatbot in the future.
    \item \textit{Effectiveness of Chatbot Interactions on Customer Satisfaction.} The research aims to determine how the effectiveness of chatbot responses, especially in resolving customer issues, affects overall satisfaction by analysing customer experiences with chatbots.
    \item \textit{The Effect of a Chatbot's Empathy and Tone on Customer Trust.} The research will explore whether using empathetic language in chatbot interactions enhances customer trust, particularly in handling sensitive inquiries.
    \item \textit{The Relationship Between Company Expectations and Customer Perceptions of Customer Service Chatbots.} The research will investigate whether there is a relationship between the benefits companies expect from implementing chatbots and the actual customer perceptions of these chatbots.
    \item \textit{The Relationship Between Customer Satisfaction and Chatbot Interactions and the Business Benefits of Chatbot Implementation.} The overarching goal is to determine how chatbot interactions influence customer satisfaction and the potential benefits for businesses in terms of customer experience.
\end{enumerate}
This research will help guide stakeholders in optimising chatbot implementation for improved customer trust, satisfaction, and an alleviated customer experience. 

\subsection{Methodology}
For the research, data will be collected using two methods. For the quantitative method it will 
focus on measuring customer perceptions and satisfaction using chatbots. Data will be collected through a Likert-scale survey. The survey will be distributed online to customers who have interacted with chatbots before.
\\
For qualitative data collection, interviews will be conducted with customers and business representatives. Participants will be selected by reaching out to businesses and customers for participation. Content and thematic analysis will be used as the technique for qualitative data analysis, focusing on identifying key patterns.

\subsection{Key Findings}
The main research question was split into four sub-research questions mentioned in the project objective. Reliability is an important factor to influence customer satisfaction, loyalty, and trust. However, reliability alone does not influence satisfaction, since interaction, overall implementation, empathy, and personalisation also influence customer satisfaction. The empathy of a chatbot improves comfort and confidence, which eventually improves customer happiness and engagement. Customers often prefer human assistance for complex or sensitive issues, but a solution to this could be enhancing the understanding and personalisation, which eventually improves customer satisfaction and effectiveness. Overall implementation and interaction of a chatbot could influence loyalty and interest.

\subsection{Conclusion}
The relationship between customer satisfaction and chatbot interaction is indicated to be present in these results. Benefits from chatbot implementation include the improvement of customer happiness, engagement, loyalty, interest, and satisfaction.

\subsection{Recommendations}
Recommendations can be provided with this conclusion, which is that companies should evaluate the implementation methods mentioned in this study. It is recommended to include empathy, reliability, and high quality interaction when implementing a chatbot.

\section{Introduction to Chatbots}

Artificial intelligence (AI) is a recent technological advancement that influenced our engagement in daily activities. Chatbots are a popular type of AI program that uses Natural Language Processing to communicate with an understandable human language. They also have a use in a business setting; some examples include customer service, e-commerce, and entertainment. They provide productivity to the user but may also entertain them or provide interaction with the company. Chatbot usage may provide some benefits for a company, like decreased labour costs or providing personalised recommendations.

\subsection{Brief History of Chatbots}

According to an article on chatbot history, the initial idea of chatbots came from Alan Turing in 1950 (Adamopoulou and Moussiades, 2020)\cite{chatbothistory}. He thought about the possibility of a computer speaking to people without them realising it. This article also stated that chatbot development started with ELIZA in 1966, which was developed to imitate a psychotherapist by rephrasing a user's answers. ELIZA had limited knowledge and did not learn from any conversations, it also cannot hold long conversations. After ELIZA, new advancements continued to be made with PARRY in 1972, the first AI chatbot Jabberwacky in 1988, TINYMUD in 1991, and eventually the first online chatbot ALICE. Eventually, voice assistants were created, like Siri in 2010. Voice assistants often require an internet connection and create responses that way. There are multiple types of chatbot approaches, some examples contain Machine Learning for AI chatbots and pattern matching chatbots, which uses rules to generate a response.
 
\subsection{Benefits of Chatbots}

In the field of customer satisfaction, the usage of chatbots is rising as they offer various values and benefits to customers. Chatbots provide fast, 24/7 customer service, as customers are willing to have solutions in a short time and at any moment. They also offer a more personalised experience, as customers understand that chatbots can collect their data solely to create a better customer experience. Additionally, chatbots provide multilingual support, catering to a global customer base and enhancing the overall user experience.
\\
Nila Armelia Windasari, an Assistant Professor and Principal Researcher at iDNA Solutions, and Moch Akbar Selamat, an Onboarding Account Manager for Google Ads at TDCX, both claim that four chatbot features are deemed important for SMEs: responsive chatbots, simple action steps to trigger customer actions, humanised conversations, and personalised recommendations. They mention that anthropomorphic cues in SMEs' chatbots may influence customer shopping intentions by creating enjoyable and useful features, which encourage customers to engage and shop with SMEs (Selamat and Windasari, 2021) \cite{Selamat2021}.



\subsection{Risks of Chatbots}

Using a chatbot may not always lead to only benefits, there will always be some risks and threats involved too.
Threats are often one-off events such as malware attacks, phishing emails, ransomware, or distributed denial of service, and also the possibility of cybercriminals threatening to expose customer data. System vulnerabilities are weaknesses that can be exploited by an attacker to cross privilege boundaries (perform unauthorised actions). The system is vulnerable when it has weak coding, or has a weak firewall, and so forth. There is an implementation to prevent this kind of risk from happening.
Authentication and authorisation: A typical example is chatting with a banking chatbot about an account balance. In this case, authentication and authorisation are necessary to identify that a user is verified with valid and secure login credentials. 
End-to-End Encryption: End-to-End Encryption is a system of communication where only the communicating parties can read the messages. The conversation is encrypted in a way that only the unique recipient of a message is allowed to decrypt it. (Hasal et al., 2021) \cite{chatbots}


\subsection{Purpose and Importance of the Research}

With the increasing use of chatbots in many businesses, especially in customer care, e-commerce, and other service sectors, comprehending their impact on consumer experiences has become valuable. Chatbots are anticipated to deliver efficiency and individualised service, thus elevating the entire user experience. Regardless, the equilibrium between these advantages and possible hazards-such as data security, absence of empathy, or even chatbot malfunction-remains underexplored. This research seeks to address the gap by assessing the efficiency of chatbots. This study will elucidate the influence of chatbot characteristics, including anthropomorphised cues, on user comfort and trust, offering essential insights for organisations to enhance chatbot design for operational efficiency and customer pleasure. (Haugeland et al., 2022) \cite {chatbotimpact}

\section{Approaches and Results of Research}

This study on a chatbot's influence on customer experience and business benefits has been conducted by students from Breda University of Applied Sciences. To answer the research question, four sub-research questions had to be answered, for which each student performed analysis individually. Two research methods were chosen. A survey with questions was shared to be used for analysing how a chatbot's empathy and tone influence the customer's trust in handling sensitive inquiries and the impact of the reliability of information provided by a chatbot on customer satisfaction. The survey's questions asked for the participant to answer with a score within a range of 1 to 5, where 1 meant that they strongly disagreed and 5 meant that they strongly agreed with the statements. Interviews with both companies and customers were also conducted to answer open-ended questions related to how customer experience influences overall satisfaction with chatbot interactions and the relationship between company expectations of benefits and customer perceptions regarding customer service. 
\\
Both analysis methods also ranged between the different research questions. These methods will be briefly explained in clear and non-technical ways below.
\\
The interviews were analysed through their specific theme, which is called thematic analysis. In these types of analyses, the number of times a specific theme was mentioned throughout interviews is evaluated. These themes and statements are afterwards interpreted to understand the participant's reasoning and eventually used to answer the research question. 
\\
The survey's data was analysed using coding techniques in the programming language Python. In this analysis, statistical methods are used to understand the relationship between two variables. These results are later interpreted and explained to ensure a clear answer to the question.

\subsection{Results Impact Perceived Reliability on Satisfaction}

An analysis of the quantitative data from the survey reveals that perceived reliability plays a crucial role in shaping customer relationships. Participants demonstrated a strong tendency to trust products and services they viewed as reliable, indicating that enhancing perceived reliability can significantly bolster customer trust. Furthermore, while perceived reliability positively influences satisfaction, the relationship is not as strong, suggesting that other factors also contribute to satisfaction levels. Most notably, perceived reliability strongly correlates with customer loyalty, as respondents who perceive high reliability are much more likely to remain loyal to a brand. These findings highlight the importance of improving perceived reliability to enhance trust and loyalty, while also considering additional factors to boost overall customer satisfaction.

\subsection{Results Influence of Chatbot Empathy on Trust}

An analysis of the quantitative data from the survey answers reveals that empathy in chatbots significantly enhances the user experience. Participants typically indicated increased comfort and confidence when engaging with empathic chatbots, particularly in sensitive domains such as customer service and sensitive inquiries. This indicates that chatbot designs including empathic replies are more likely to improve customer happiness and engagement, enabling businesses to deliver more personalised and helpful experiences without compromising efficiency.

\subsection{Results Expected Benefits and Customer Perception}

The analysis on the sub-research question: 'Is there a relationship between the company expectations regarding the benefits of customer service chatbots and the customer perceptions of their experiences?' was conducted using both qualitative and quantitative methods. Interviews were used to answer the main question. This analysis suggests that there is evidence of a relationship between the two, as both statements from the companies and customers align well. This was concluded by comparing both the customer and company ideas of the benefits and their experiences. A common theme was how the chatbot interaction quality has an impact on the customer's loyalty towards a company. This was further evaluated through analysing survey questions, which also indicated a strong relationship between interaction and loyalty.

\subsection{Results Relationship Between Chatbot Effectiveness and Customer Satisfaction}

Chatbot performance expectations were a common theme in both qualitative and quantitative analysis. Because of predetermined and generic responses, informants disliked chatbots. Customers are frustrated by the chatbot's inability to provide personalised solutions for complex or personal issues. While some customers understood chatbot basics, their expectations varied based on prior experience. Those less familiar with chatbot capabilities had unrealistic expectations for handling complex queries, while those more familiar acknowledged its technological limitations but were disappointed in its overall effectiveness.
For sensitive topics requiring personalised support, customers preferred human customer service representatives. Interviewees suggested chatbots provide direct, reliable answers without unnecessary or excessive answers when accurate. To complete that, chatbots that learn from past interactions and have personalised responses were suggested to improve effectiveness.
Numerous participants acknowledged chatbots' effectiveness, but opinions differ on personalisation and problem-solving. Users who found chatbots easy to use were more likely to believe they could solve their problems. Chatbots still fail to meet user expectations, especially in complex situations where human representatives are preferred.
Improved contextual understanding of chatbots could boost user satisfaction and effectiveness. Enhancing chatbot capabilities to provide more personalised and context-accurate answers could bridge the gap between customer expectations and actual performance, ultimately leading to higher satisfaction.

\section{Conclusions}

By interpreting all sub-research questions, the answer to the main research question can be concluded: "What is the relationship between customer satisfaction and chatbot interaction and what specific benefits do they offer to companies in terms of customer experience?" First of all, reliability is an important factor that influences customer satisfaction, loyalty, and trust. Reliability alone does not influence satisfaction, as it is suggested that other factors also contribute to satisfaction. As proven by the results of the research on expected benefits, these other factors contain interaction and implementation, which will eventually influence multiple benefits. These benefits contain increased loyalty and interest, which were proven to be related. The research on chatbot empathy has also indicated a relationship between interaction and satisfaction. That research revealed that engaging with empathetic chatbots increased comfort and confidence, improving customer happiness and engagement. To further support this, the research on the relationship between chatbot effectiveness and customer satisfaction revealed that customers' expectations often fell short for complex or sensitive issues, where customers often prefer human assistance. By enhancing understanding and personalisation, this gap could be bridged, which could eventually improve customer satisfaction and effectiveness.
\\
In conclusion, the relationship between customer satisfaction and chatbot interaction is indicated to be present when combining the results of all research questions. These two features cause a positive relationship, where the increase in chatbot interaction quality also increases customer satisfaction. Important benefits of chatbot implementation are the improvement of customer happiness, engagement, loyalty, interest, and satisfaction. These benefits are provided through a chatbot's empathy, overall implementation, reliability, understanding, personalisation, and interaction quality. 

\section{Recommendations}

After investigating the relationship between customer satisfaction and chatbot interaction, recommendations could be made for future chatbot implementation and satisfaction management. When implementing a customer service chatbot in the future, it is recommended to evaluate implementation methods and ideas and include empathy, reliability, and high quality interaction. Some of these features, such as empathy and the implementation methods, would require evaluation based on the company's needs to influence actual benefits.

\begin{thebibliography}{99}

\bibitem{chatbothistory}
Adamopoulou, E., and Moussiades, L. (2020). Chatbots: History, technology, and applications. Machine Learning with Applications, 2, 100006. \url{https://doi.org/10.1016/j.mlwa.2020.100006}

\bibitem{securityrisks}
Edu, J., Mulligan, C., Pierazzi, F., Polakis, J., Suarez-Tangil, G., and Such, J. (2022). Exploring the security and privacy risks of chatbots in messaging services. ACM Digital Library. \url{https://doi.org/10.1145/3517745.3561433}


\bibitem{Selamat2021}
Selamat, Moch Akbar, and Nila Armelia Windasari. 2021. "Chatbot for SMEs: Integrating Customer and Business Owner Perspectives." \textit{Technology in Society} 66: 101685. \url{https://doi.org/10.1016/j.techsoc.2021.101685.}

\bibitem{chatbotimpact}
Haugeland, I. K. F., Følstad, A., Taylor, C., and Bjørkli, C. A. (2022). Understanding the user experience of customer service chatbots: An experimental study of chatbot interaction design. \url{https://www.sciencedirect.com/science/article/pii/S1071581922000179}

\bibitem{chatbots}
Hasal, M., Nowaková, J., Saghair, K. A., Abdulla, H., Snášel, V., and Ogiela, L. (2021). Chatbots: Security, privacy, data protection, and social aspects. Concurrency and Computation: Practice and Experience. \url{https://doi.org/10.1002/cpe.6426}



\end{thebibliography}

\section*{Appendix 1}
\textbf{Stakeholder Categories:}
\begin{itemize}
 \item \textbf{A.} Hight attention, High power
 \item \textbf{B.} Low attention, High power
 \item \textbf{C.} High attention, Low power
 \item \textbf{D.} Low attention, Low power
\end{itemize}
Below, all our important stakeholders are sorted by their category:
\begin{itemize}
 \item \textbf{A.} Small businesses, Client: Digiwerkplaats, Developers
 \item \textbf{B.} Industry Regulators
 \item \textbf{C.} Customers, Customer Support Teams
 \item \textbf{D.} Not applicable
\end{itemize}
Of these, there are also Main stakeholders, which have specific needs and requirements:
\begin{itemize}
    \item \textbf{Main Stakeholders:} Small businesses and our client
    \item \textbf{Needs:} To understand the impact of chatbots on customer attitudes and how beneficial they could be to the company.
\end{itemize}

\textbf{Impact of Research:}
\begin{itemize}
    \item \textbf{Small Businesses and client:} The small businesses and clients are the main interest of our research. They impact our research by giving insights about the possible benefits that chatbots could have.
    \item \textbf{Developers:} Developers are important for the research since they are the ones to eventually get interested in the research and use gained knowledge on the development of chatbots. A developer could impact the research by previously created chatbots and the chatbot’s User Experience.
    \item \textbf{Industry Regulators:} Regulation is very important for both chatbots and research papers. For industry Rrgulators, it would be very important for our research to follow their regulations. Not following these guidelines and regulations correctly might cause legal or ethical problems.
    \item \textbf{Customers:} Customers will be part of the people to be interviewed, and they would impact the research by providing data related to their opinions and experiences with chatbots.
    \item \textbf{Customer Support Teams:} Customer support teams are the ones to likely work closely with chatbots, and for them, it is important to know how customers react towards chatbots and what to change for them. Customer Support Teams might also provide meaningful insights about chatbot usage in customer service scenarios.
\end{itemize}

\end{document}